\subsection{\texorpdfstring{SUBGROUPS AND QUOTIENT GROUPS OF \(T^{n}\)}{SUBGROUPS AND QUOTIENT GROUPS OF T TO THE N}}

Let us identify \(\mathbf{T}^n\) with \(\mathbf{R}^n/\mathbf{Z}^n\) , and let \(\varphi\) be the canonical homomorphism of \(\mathbf{R}^n\) onto \(\mathbf{R}^n/\mathbf{Z}^n\) . Every subgroup of \(\mathbf{T}^n\) is of the form \(\mathbf{G} = \varphi(\mathbf{H})\) where \(\mathbf{H}\) is a subgroup of \(\mathbf{R}^n\) which contains \(\mathbf{Z}^n\) and is isomorphic to \(\mathbf{H}/\mathbf{Z}^n\) (Chapter III, § 2, no. 7, Proposition 20); for \(\mathbf{G}\) to be closed in \(\mathbf{T}^n\) it is necessary and sufficient for \(\mathbf{H}\) to be closed in \(\mathbf{R}^n\) (Chapter I, § 3, no. 4). Thus in order to find the closed subgroups of \(\mathbf{T}^n\) we have to determine all the closed subgroups \(\mathbf{H}\) of \(\mathbf{R}^n\) which contain \(\mathbf{Z}^n\) ; to do this we shall use Proposition 6 of no. 3, and we start by determining the subgroup \(\mathbf{H}^*\) associated with \(\mathbf{H}\) . Since \(\mathbf{Z}^n\) is associated with itself, we have \(\mathbf{H}^* \subset \mathbf{Z}^n\) ; consequently (no. 1) there exists a basis \((\boldsymbol{a}_i)_{1 \leqslant i \leqslant n}\) of \(\mathbf{R}^n\) which generates \(\mathbf{Z}^n\) , and a basis of \(\mathbf{H}^*\) (with respect to the ring \(\mathbf{Z}\) ) consisting of \(p\) points \(\boldsymbol{b}_i\) ( \(1 \leqslant i \leqslant p\) ) such that \(\boldsymbol{b}_i = e_i \boldsymbol{a}_i\) ( \(1 \leqslant i \leqslant p\) ), where the \(e_i\) are integers satisfying the congruences \(e_{i+1} \equiv o (\text{mod } e_i)\) for \(1 \leqslant i \leqslant p - 1\) . Let \((\boldsymbol{a}_i')\) be the basis dual to \((\boldsymbol{a}_i)\) ; then \(\boldsymbol{u} = \sum_{i=1}^{n} u_i \boldsymbol{a}_i'\) belongs to \((\mathrm{H}^*)^* = \mathrm{H}\) if and only if \(u_i e_i \in \mathbf{Z}\) for \(1 \leqslant i \leqslant p\) ; in other words, \(\mathbf{H}\) is the direct sum of the vector subspace V generated by \(\boldsymbol{a}_{p+1}'\) , \ldots, \(\boldsymbol{a}_n'\) , and the discrete subgroup K generated by the \(p\) points

\[
e _ {i} ^ {- 1} \boldsymbol {a} _ {i} ^ {\prime} \quad (\mathrm{I} \leqslant i \leqslant p).
\]

On the other hand, \(\mathbf{Z}^n\) is the direct sum of \(\mathbf{V} \cap \mathbf{Z}^n\) and \(\mathbf{K} \cap \mathbf{Z}^n\) , because the \(a_i'\) ( \(i \leqslant i \leqslant n\) ) generate \(\mathbf{Z}^n\) . The quotient group \(\mathrm{H} / \mathbf{Z}^n\) is therefore isomorphic to \((\mathrm{V} / (\mathrm{V} \cap \mathbf{Z}^n)) \times (\mathrm{K} / (\mathrm{K} \cap \mathbf{Z}^n))\) (Chapter III, § 2, no. 9, Corollary to Proposition 26); \(\mathrm{V} / (\mathrm{V} \cap \mathbf{Z}^n)\) is isomorphic to \(\mathbf{T}^{n-p}\) , and \(\mathbf{K} / (\mathbf{K} \cap \mathbf{Z}^n)\) is a finite group, the direct sum of \(p\) cyclic groups of orders \(e_i\) , respectively ( \(i \leqslant i \leqslant p\) ) (cf.~no. 1).

Using the same notation, every Hausdorff quotient group of \(T^{n}\) is of the form \(\mathbf{T}^{n}/\varphi(\mathbf{H})\) and is isomorphic to \(R^{n}/H\) (Chapter III, § 2, no. 7, Corollary to Proposition 22); if W is the vector subspace generated by K, \(R^{n}/H\) is isomorphic to W/K (Chapter III, § 2, no. 9, Corollary to Proposition 26), i.e., to \(T^{p}\) . To sum up:

PROPOSITION 11. Every closed subgroup of \(\mathbf{T}^n\) is isomorphic to a group of the form \(\mathbf{T}^h \times \mathbf{F}\) ( \(0 \leqslant h \leqslant n\) ) where \(\mathbf{F}\) is a finite abelian group, such that the smallest number of cyclic subgroups of which \(\mathbf{F}\) is a direct sum is \(\leqslant n - h\) . Every Hausdorff quotient group of \(\mathbf{T}^n\) is isomorphic to a group of the form \(\mathbf{T}^h\) ( \(0 \leqslant h \leqslant n\) ).

In particular \((n = 1)\) :

COROLLARY. Every closed subgroup of T, other than T itself, is a finite cyclic group. Every Hausdorff quotient group of T, other than \(\{o\}\) , is isomorphic to T.

